\section{Conclusion}
\label{sec:conclusion}

This systematic review analyzed 13 studies on wearable seizure detection and forecasting using deep learning approaches. The evidence base comprises 912 participants across nine detection studies and four forecasting studies published between 2013 and 2026.

The predominance of accelerometer-based detection reveals a fundamental constraint in wearable seizure monitoring. Six of nine detection studies rely on accelerometry, and the two studies achieving ILAE Phase 3 benchmarks for false positive rates both use ACC. \textcite{spahrDeepLearningbasedDetection2025} achieved 0.0054/h FPR with single-modality ACC. \textcite{fineDetectionKeyAutomated2025} achieved 0.023/h FPR with a 6-axis ACC and gyroscope band. This convergence on a single modality reflects the biological reality that convulsive seizures produce distinct motor signatures that deep learning can extract reliably. However, it also means current wearable approaches can serve only the subset of patients with motor manifestations. Approximately 30\% of people with epilepsy have non-convulsive seizures that remain undetectable with movement-based approaches.

Forecasting studies face a different constraint. The consistent AUC plateau around 0.74 to 0.75 across studies using different methods, modalities, and patient populations suggests a fundamental limit on pre-ictal predictability using non-invasive sensors. Unlike detection, where the ictal motor event produces a clear signal, forecasting must rely on autonomic precursors that may be inherently less specific. The autonomic nervous system responds to stress, exercise, sleep transitions, and other daily events. These confounds may establish a performance ceiling for wearable forecasting that cannot be breached without EEG biomarkers or novel sensing modalities.

No study in this review has achieved all requirements for widespread clinical adoption. Commercial detection devices exist with FDA or CE approval, but the deep learning approaches reviewed here lack large-scale prospective validation. Forecasting faces an additional barrier. No device has regulatory approval for forecasting, and the field lacks established clinical benchmarks equivalent to the ILAE Phase 3 FPR standard. The performance achieved to date, while scientifically promising, may not be sufficient for regulatory approval or clinical utility.

The clinical pathway forward requires addressing three gaps. First, home validation at scale is needed to bridge the EMU-to-home translation gap. The two studies meeting FPR benchmarks were conducted in hospital settings. Real-world deployment introduces activities and confounds not captured in EMU environments. Second, forecasting needs standardized clinical benchmarks to define what level of AUC, IoC, or TiW constitutes clinically meaningful performance. Third, novel modalities or approaches are needed for non-convulsive seizures. ECG-based detection shows some promise for identifying autonomic changes during non-motor seizures, but current performance remains insufficient. Mixed personalization strategies that combine population patterns with individual adaptation may address the deployability versus performance trade-off without requiring months of individual data collection.
